\section{Conclusion}

A probability returned with every typed judgment should decide when a second
judgment is worth buying. Our experiments on frozen Laya and Vela across three
tasks do not support that promise: evolutionary search produces no held-out
advantage over fixed policies or budget-matched random search.

Gross complementarity is real---a per-row oracle gains up to 23.3 accuracy
points---but harmful switches offset repairs, and the net balance is small,
metric-dependent and unstable across pools. Calibration and action-value
prediction dissociate: temperature scaling
improves ECE and Brier but not row order, and in both FinDVer cells the
probability ranks repairable rows below chance. The search landscape is flat:
most mutations are neutral, most runs plateau at generation zero, and the
training ranking does not predict the dev ranking.

Limits are part of the result: pools are small, the \noul{} template was not
audited, only two open models were run, and the locked test remains unrun. Next:
a supervision target judged by held-out net utility.